\begin{lemma}[Situation 7.2, Subcase 2.1 Gamma exception]
Let \(D\) be a \noderef{HierarchySituationTwoCaseTwoData}.  Suppose
there is a color \(\eta_h\) such that, at the last divider index \(q\),
\[
 \eta_h\in D_{n,q},\qquad
 \eta_h\in\overline\varphi_n(y_{p-1}),\qquad
 \theta\in\Gamma^q_{\eta_h}.                               \tag{1}
\]
Then a contradiction follows.
\end{lemma}
\begin{proof}
Put \(H=C.\mathit{hierarchy}\),
\(A=\noderef{HierarchyStarredStage}(H,q)\),
\(\ell=\noderef{HierarchyFullTailLength}(C)\), and
\[
 B=\noderef{HierarchyFullTailPrefix}(C,\ell-2),\qquad
 B^+=\noderef{HierarchyFullTailPrefix}(C,\ell-1).
\]
As in the exact prefix definitions,
\[
 A\preccurlyeq B\preccurlyeq B^+\preccurlyeq T,\qquad
 |B|=|T|-2,\quad |B^+|=|T|-1.                              \tag{2}
\]
Apply \noderef{HierarchyCounterexamplePenultimateElementary} to the six
fields of the supplied \noderef{HierarchyMinimumCounterexampleContext};
then
\[
              V(B^+)\text{ is elementary under }\varphi_n. \tag{3}
\]
Apply \noderef{HierarchySituationTwoCaseTwoExclusions} as well.  It gives
\[
 \alpha\notin\overline\varphi_n(B),\qquad
 \alpha\notin D_n,\qquad
 \alpha\notin\varphi_n\langle T-A\rangle.                  \tag{4}
\]

By the two-color field of \noderef{GoodHierarchy}, the Gamma set in
(1) has exactly two members.  Since it contains \(\theta\), choose its
other member \(\gamma_2\); explicitly,
\[
 \Gamma^q_{\eta_h}=\{\theta,\gamma_2\},\qquad
 \theta\ne\gamma_2.                                       \tag{5}
\]
Membership of \(\eta_h\) in
\noderef{HierarchyStoredColorsAt} says it is not missing on \(A\).
It is missing at \(y_{p-1}\), and (3) shows it cannot also be missing at
any earlier vertex of \(B\).  Thus \(y_{p-1}\) is its first missing
vertex.  The avoidance field of \noderef{GoodHierarchy}, applied to
(5), gives
\[
 \gamma_2\in\overline\varphi_n(A),\qquad
 \gamma_2\notin\varphi_n\langle B^+-A\rangle.               \tag{6}
\]
The last edge has color \(\theta\ne\gamma_2\), so (6) excludes
\(\gamma_2\) from every exact edge of \(T-A\).  Also
\(\gamma_2\ne\alpha\), since \(\gamma_2\) is missing on \(A\) whereas
(4) implies that \(\alpha\) is not missing there.

Apply \noderef{ChangeLemma} with ordered colors
\((\gamma_2,\alpha)\).  Both colors are missing on \(B^+\), witnessed
in \(A\) and at \(y_{p-1}\), respectively.  For positive divider count,
the first color is missing on the last ordinary stage; at divider zero it
is missing on \(T_n\), exactly as supplied by the stage-missing part of
avoidance.  Because \(\alpha\) is not missing on \(A\), the conditional
order premise is active: the first \(\gamma_2\)-missing vertex
\(v_{\gamma_2}\in V(A)\) precedes \(y_{p-1}\), and (6) is the exact
unused interval.  The other inputs are (3) and the six minimum-context
fields.  Hence
\[
 K_{\alpha,\gamma_2}(v_{\gamma_2})
       =K_{\alpha,\gamma_2}(y_{p-1}).                        \tag{7}
\]

This component is disjoint from
\(Q_1=K_{\alpha,\gamma_2}(y_p)\).  Indeed,
\(v_{\gamma_2},y_{p-1},y_p\) are distinct by (2) and vertex injectivity
in \noderef{TreeSequence}; the first misses \(\gamma_2\), while the
other two miss \(\alpha\).  If the components were equal, properness,
\noderef{TwoColorIncidenceDegreeAtMostTwo}, and
\noderef{TwoColorComponentSpanningPath} would put three distinct
degree-at-most-one vertex-nodes at the two endpoints of one finite path.
Thus \noderef{TwoColorComponentsDisjoint} gives disjointness, and the
spanning-path result supplies the required
\noderef{TwoColorPathComponent} certificate, including the isolated
vertex possibility.

Define the first recoloring explicitly by
\[
 \mu_1=\noderef{TwoColorComponentSwap}
       (G,\varphi_n,\gamma_2,\alpha,\operatorname{inl}y_p). \tag{8}
\]
Apply \noderef{StableChangeCasesOneTwo} before using any stability
conclusion.  Use the ordered colors \((\gamma_2,\alpha)\), (3), (6)--(8),
the stage-missing alternative just used for the Change Lemma, and all six
minimum-context inputs.  Retain in particular its
\noderef{HierarchyWeaklyStable} conclusion for \(\mu_1\).

We now verify, independently, every premise needed by the full-tail
provider.  First fix \(u\in V(B^+)\).  If its vertex-node is outside
\(Q_1\), no incident edge is changed.  If it lies in \(Q_1\), it cannot
be an endpoint, since an endpoint misses one swapped color and (3), (7),
and component disjointness put the unique such witness in the other
component.  It is therefore an internal vertex.  Properness gives exactly
one incident \(\alpha\)-edge and one incident \(\gamma_2\)-edge, so the
swap exchanges them and preserves the full set of missing colors.  Thus
\[
 \overline\mu_1(u)=\overline\varphi_n(u)
                    \quad(u\in V(B^+)).                     \tag{9}
\]

Next fix an exact tail-edge position.  Equation (4) excludes \(\alpha\)
and (6), together with the last-edge equation and (5), excludes
\(\gamma_2\).  Consequently its individual edge-node is outside the
swapped component and its value is fixed.  This check includes the
possible divider-zero SE/RE seed: if one instead reaches that edge through
the exceptional conclusion of
\noderef{HierarchyTailMissingColorExclusion}, the weak stability obtained
before this paragraph is a \noderef{StableColoring} relation on \(T_n\); the exact
singleton connecting-edge certificate makes the seed incident with
\(T_n\), so the controlled-edge clause fixes that individual edge.  Hence
\[
 \mu_1(T.\mathit{edge}(t))=\varphi_n(T.\mathit{edge}(t))
                    \quad(|A|\le t<|T|).                    \tag{10}
\]
No equality follows merely from component separation.

Finally consider the least missing position of an arbitrary stored color.
The color \(\gamma_2\) is not stored because it is missing on \(A\), and
\(\alpha\) is not even in the raw set \(D_n\), by (4).  Every stored
color is therefore different from both swapped colors, so its missing
predicate is pointwise unchanged at every vertex; unfolding
\noderef{HierarchyFirstMissingPosition} gives equality of its totalized
least positions.  Thus (9), (10), and the stored-position equality verify
the three separate transport premises of
\noderef{StableChangeCasesOneTwoFullTail}.  Applying that theorem yields
an actual corresponding ETT on the exact tree \(T\), maximum property, an
actual base, a full good hierarchy, and
\noderef{SameGoodHierarchy} under \(\mu_1\).

We package these returned objects as a genuine
\noderef{HierarchyReplacementPresentation} \(P_1\).  The component swap
is proper by \noderef{TwoColorComponentSwapProper} and fixes the exact
uncolored-edge set, giving the properness and
\noderef{IsOneEdgeUncolored} fields.  The swap makes \(\gamma_2\) missing
at \(y_p\), while component disjointness retains its missing witness
\(v_{\gamma_2}\); these distinct witnesses make the returned good
hierarchy a \noderef{HierarchyCounterexample}.  The corresponding ETT has
the old rung.  Copy the complete rung-uniform lower assertion.  Same-good
gives the old divider count, so the old fewer-divider field applies.
Because the terminal tree is literally \(T\) and same-good identifies the
ordinary stages, both coordinates of
\noderef{HierarchyCounterexampleMeasure} are equal to the old ones; this
transports lexicographic minimality.  Weak stability identifies missing
sets on the terminal ladder tree and transports the \(2n+11\) bound.
These arguments fill, separately, all six fields of the replacement
\noderef{HierarchyMinimumCounterexampleContext}.  Correspondence and
same-good are composed with the original presentation entry by entry,
using the transitivity statement in \noderef{sc2} for every positive
stable-coloring slot.

Suppose first that \(\eta_h\) is missing at \(y_p\) under \(\mu_1\).
Construct a \noderef{HierarchySituationTwoCaseOneData} for \(P_1\) as
follows.  Use the original exact edge and vertex indices and non-incidence
equation.  Literal equality of \(T\) and same-good preserve the full-tail
length and path number.  Equations (9)--(10) preserve the last-edge color
\(\theta\) and its presence at \(y_{p-1}\).  Use \(\eta_h\) as the Case
1 color; it is missing at both terminal vertices by hypothesis and (1),
and same-good transports its stored membership and the pair
\(\{\theta,\gamma_2\}\) from (5).  The two members are distinct.  Every
field is now present, so
\noderef{HierarchySituationTwoSubcaseOneOne}, applied to this actual Case
1 datum and the actual context of \(P_1\), gives a contradiction.

It remains to assume that \(\eta_h\) is present at \(y_p\) under
\(\mu_1\).  Under \(\mu_1\), the color \(\gamma_2\) is missing at both
\(v_{\gamma_2}\) and \(y_p\), while \(\eta_h\) remains missing at
\(y_{p-1}\) by (9).  Its stored membership says it is not missing on
\(A\), and elementarity transported by (9) says its first missing vertex
is \(y_{p-1}\).  Equation (10) and (6) show that \(\gamma_2\) is unused
on the full tail under \(\mu_1\).  The colors are distinct because
\(\gamma_2\) is missing on \(A\) while \(\eta_h\) is not.  Moreover,
same-good transports (5), and avoidance in the replacement
\noderef{GoodHierarchy} again places \(\gamma_2\) in the required
replacement stage-missing set.  Apply \noderef{ChangeLemma} inside
\(P_1\), with ordered colors \((\gamma_2,\eta_h)\), the explicit first
positions, this exact unused premise, the stage-missing alternative for
\(\gamma_2\), and all fields of the transported minimum context.  It
gives
\[
 K_{\eta_h,\gamma_2}^{\mu_1}(v_{\gamma_2})
       =K_{\eta_h,\gamma_2}^{\mu_1}(y_{p-1}).                \tag{11}
\]
The component in (11) is disjoint from the component rooted at \(y_p\):
otherwise the three distinct vertices \(v_{\gamma_2},y_{p-1},y_p\),
missing respectively \(\gamma_2,\eta_h,\gamma_2\), would again be three
endpoints of the spanning path.  Let
\[
 \mu_2=\noderef{TwoColorComponentSwap}
       (G,\mu_1,\gamma_2,\eta_h,\operatorname{inl}y_p).      \tag{12}
\]
Apply \noderef{StableChangeCasesOneTwo} immediately, relative to \(P_1\),
before using stability to protect any edge.  The stage-missing color is
again \(\gamma_2\); equation (11), the disjointness just proved, and the
full-tail unusedness verify its remaining Cases 1/2 premises.

We again check the full-tail inputs one by one.  The endpoint/internal
vertex argument used for (9), now with (11), proves
\[
 \overline\mu_2(u)=\overline\mu_1(u)
                    \quad(u\in V(B^+)).                     \tag{13}
\]
The color \(\gamma_2\) is unused on every tail edge by (6) and (10).
Apply \noderef{HierarchyTailMissingColorExclusion} to \(\eta_h\), which
is missing at the last vertex of the elementary tree \(B^+\).  It excludes
\(\eta_h\) from every exact tail edge except possibly the individual
divider-zero SE/RE seed carrying a raw stored color.  The weak stability
obtained before this edge argument is a \noderef{StableColoring} relation
on \(T_n\) and fixes precisely that seed through its
controlled incident-edge clause.  The last edge has color \(\theta\),
which is distinct from \(\eta_h\) because the former is present and the
latter missing at \(y_{p-1}\), and distinct from \(\gamma_2\) by (5).
Therefore every exact tail-edge value is fixed:
\[
 \mu_2(T.\mathit{edge}(t))=\mu_1(T.\mathit{edge}(t))
                    \quad(|A|\le t<|T|).                    \tag{14}
\]

For stored least positions, \(\gamma_2\) is not stored.  The affected
stored color \(\eta_h\) has first missing position \(y_{p-1}\), as proved
above, and (13) identifies every candidate through that witness; hence its
least position is unchanged.  Every other stored color is unaffected.
Thus the full quantified stored-position equality also holds.  Apply
\noderef{StableChangeCasesOneTwoFullTail} with (13), (14), and this
equality.  It supplies an actual corresponding same-rung maximum ETT,
base, hierarchy, goodness, and same-good data under \(\mu_2\).

Build \(P_2\) field by field exactly as for \(P_1\).  Properness and the
exact uncolored-edge property come from the swap theorem and its
definition.  The swap makes \(\eta_h\) missing at \(y_p\), while (13)
retains its missingness at the distinct vertex \(y_{p-1}\), so these
witnesses provide non-elementarity.  Literal equality of \(T\), rung
agreement, and same-good transport the lower-rung field, the
fewer-divider field, equality of both measure coordinates and hence
minimality, and the large-set bound, separately.  Compose the two
correspondence certificates, every positive-stage stable relation, and
the two same-good relations to obtain an original-relative
\noderef{HierarchyReplacementPresentation} with all six context fields.

Finally construct a \noderef{HierarchySituationTwoCaseOneData} for
\(P_2\).  Copy the original exact tail indices, non-incidence, full-tail
length, and path number.  Equations (10), (13), and (14) preserve the last
edge's color \(\theta\) and its presence at \(y_{p-1}\).  Take
\(\eta_h\) as the Case 1 color; it is missing at both terminal vertices,
its stored membership is transported, and its Gamma pair remains exactly
\(\{\theta,\gamma_2\}\) with distinct members.  Applying only the already
proved \noderef{HierarchySituationTwoSubcaseOneOne} to this actual datum
and the actual minimum context of \(P_2\) yields the final contradiction.
\end{proof}
